\documentclass[11pt]{article}
\usepackage[a4paper,margin=1in]{geometry}
\usepackage[T1]{fontenc}
\usepackage{helvet}
\renewcommand{\familydefault}{\sfdefault}
\usepackage{xcolor}
\usepackage{booktabs}
\usepackage{array}
\usepackage{colortbl}
\usepackage[font=footnotesize,labelformat=empty]{caption}
\usepackage{titlesec}
\usepackage{fancyhdr}
\usepackage{enumitem}
\usepackage{amsmath}
\usepackage{lastpage}
\usepackage[hidelinks]{hyperref}

% ---- Bayesian Capital brand palette ----
\definecolor{bcnavy}{RGB}{11,31,58}
\definecolor{bcblue}{RGB}{32,74,135}
\definecolor{bcgold}{RGB}{176,141,87}
\definecolor{bcgrey}{RGB}{90,98,112}
\definecolor{bcrule}{RGB}{210,214,220}
\definecolor{bckeep}{RGB}{232,237,244}
\definecolor{bcact}{RGB}{247,241,228}

\titleformat{\section}{\color{bcnavy}\Large\bfseries}{\thesection}{0.6em}{}
\titleformat{\subsection}{\color{bcblue}\large\bfseries}{\thesubsection}{0.6em}{}
\titlespacing*{\section}{0pt}{1.4em}{0.5em}

\pagestyle{fancy}\fancyhf{}
\renewcommand{\headrulewidth}{0pt}
\renewcommand{\footrulewidth}{0.4pt}
\renewcommand{\footrule}{\hbox to\headwidth{\color{bcrule}\leaders\hrule height \footrulewidth\hfill}}
\fancyfoot[L]{\footnotesize\color{bcgrey}Bayesian Capital}
\fancyfoot[C]{\footnotesize\color{bcgrey}Confidential --- Internal Research}
\fancyfoot[R]{\footnotesize\color{bcgrey}Page \thepage\ of \pageref{LastPage}}

\newcommand{\kpi}[2]{\textbf{#1}\,{\footnotesize\color{bcgrey}#2}}
\setlength{\parindent}{0pt}\setlength{\parskip}{0.55em}

\begin{document}

% ---------------- Title block ----------------
\begin{titlepage}
\vspace*{1.2cm}
{\color{bcgold}\rule{\textwidth}{2pt}}\\[0.4cm]
{\color{bcnavy}\Huge\bfseries Stop-Loss Period Study}\\[0.35cm]
{\color{bcblue}\LARGE Findings Memorandum}\\[0.5cm]
{\color{bcgold}\rule{\textwidth}{1pt}}\\[0.6cm]
{\large\color{bcgrey} Hybrid Bayesian / Ornstein--Uhlenbeck Volatility Harvester}\\[0.2cm]
{\large\color{bcgrey} Ten-name high-beta US equity book --- calendar-day stop (cell T2)}\\[1.2cm]
{\color{bcnavy}\large\bfseries Bayesian Capital}\\[0.15cm]
{\color{bcgrey} Quantitative Research}\\[0.15cm]
{\color{bcgrey} 29 July 2026}\\[2cm]

\fbox{\parbox{0.92\textwidth}{\color{bcnavy}\textbf{One-line conclusion.}\\ \color{black}
The uniform 50-calendar-day stop is not optimal, and the intuition that different names
want different stops is correct --- but the effect resolves into \emph{two coarse regimes},
not ten bespoke numbers. Momentum names want a \emph{longer} leash; choppier names want a
\emph{tighter} one. The single highest-value change is \textbf{RKLB $50\rightarrow60$}, which
raises return \emph{and} Sharpe \emph{and} cuts drawdown simultaneously.}}
\end{titlepage}

% ---------------- Executive summary ----------------
\section{Executive summary}

This memo measures the impact of varying the calendar-day stop-loss period (workbook cell
\texttt{T2}) for each of the ten names, holding every other parameter frozen. All figures come
from a Python replica of the workbook that reproduces \textbf{every one of the ten Model sheets
to the dollar} at the current settings.

\begin{itemize}[leftmargin=1.4em,itemsep=2pt]
  \item \kpi{Two regimes, not ten dials.}{} For trending names the stop \emph{almost never
        fires} --- it is a disaster backstop, and shortening it merely amputates compounding
        winners. For choppier names (failed bounces) a tighter stop cuts losers sooner,
        lifting risk-adjusted return and trimming drawdown.
  \item \kpi{The broad move.}{} A \textbf{$\sim$40-day} stop weakly dominates the current 50
        across the book (better or tied on return for 8/10 names), with one clear exception.
  \item \kpi{The exception (adopt first).}{} \textbf{RKLB} wants the \emph{opposite}: lengthening
        $50\rightarrow60$ lifts annual return $587\%\rightarrow732\%$, Sharpe
        $3.55\rightarrow5.10$, and \emph{reduces} max drawdown $34.5\%\rightarrow26.7\%$. The
        50-day stop is cutting the book's most explosive winner one bounce too early.
  \item \kpi{Out-of-sample.}{} Tuning the stop on the past beats the frozen 50 out of sample in
        \textbf{6/10} names --- better than a coin flip, and far better than full parameter
        re-optimisation (which loses). A single, economically-meaningful knob generalises; a
        ten-dimensional refit does not. \emph{But} the exact optimum is not perfectly stable, so
        use round numbers in the flat zone, never the in-sample argmax.
\end{itemize}

% ---------------- Method ----------------
\section{Scope and method}

The stop-loss rule exits a held position once it has been open for \texttt{T2} calendar days
without reaching its take-profit target, selling at the prevailing open. We swept
\texttt{T2} across 10, 15, 20, 25, 30, 40, 50, 60, 75, 90, 120 and 150 days plus a no-stop
case, per stock, all other
parameters frozen at their workbook values, and recorded annualised return, Sharpe, maximum
drawdown, buy count and stop-loss exits. Data: the current ten-name workbook (Apr 2024
onward, 580 trading days, 2.2-year annualisation horizon, matching cell~Y5).

Before trusting any sweep, the engine was validated against each Model sheet's cached results.
It reproduces all ten \textbf{exactly} --- e.g.\ NVDA \$6{,}517{,}886 / 184 buys / 0 stops;
RKLB \$68{,}371{,}562 / 371 buys / 1 stop --- so every number below is a faithful re-run of the
workbook's own logic, not an approximation.

% ---------------- The mechanism ----------------
\section{Why one number cannot fit all ten}

The decisive observation is that for the strongest trending names, the sweep is
\emph{flat} from roughly 50 days upward: NVDA, VST, SPOT and RKLB return \emph{identical}
results at \texttt{T2}~$=50,75,90,\dots,\infty$. The stop simply never binds --- these names
trend through their holding windows and exit at target long before any calendar limit.
For such names the stop is insurance against a disaster that did not occur in-sample;
lengthening it changes nothing, while \emph{shortening} it forces early exits that cut
winners and lower both return and Sharpe.

The choppier names behave oppositely. TSLA, PLTR and AVGO suffer failed mean-reversions ---
positions that drift sideways or down after entry. Here a shorter fuse ($\sim$40, or even 20
for TSLA) releases trapped capital sooner, improving risk-adjusted return and cutting
drawdown. The stop is an \emph{active} risk control, not a backstop.

% ---------------- Per-name table ----------------
\section{Per-name results and recommendation}

\begin{table}[h]\centering\small
\renewcommand{\arraystretch}{1.25}
\begin{tabular}{l r r r c r r r c}
\toprule
& \multicolumn{3}{c}{\textbf{Current (T2\,=\,50)}} & & \multicolumn{3}{c}{\textbf{At suggested T2}} & \\
\cmidrule(lr){2-4}\cmidrule(lr){6-8}
\textbf{Name} & ann & Sharpe & maxDD & \textbf{T2} & ann & Sharpe & maxDD & \textbf{regime}\\
\midrule
\rowcolor{bckeep} NVDA & 150\% & 4.54 & 12.5\% & \textbf{50} & 150\% & 4.54 & 12.5\% & trender (inert)\\
TSM  & 133\% & 4.76 & 11.1\% & \textbf{40} & 134\% & 4.79 & 11.1\% & mild\\
TSLA & 157\% & 2.71 & 20.3\% & \textbf{40} & 157\% & 2.57 & 22.3\% & choppy\\
VRT  & 219\% & 4.25 & 15.4\% & \textbf{40} & 235\% & 4.51 & 16.4\% & choppy\\
\rowcolor{bckeep} VST  & 257\% & 5.19 & 15.5\% & \textbf{50} & 257\% & 5.19 & 15.5\% & trender\\
AVGO & 142\% & 4.31 & 17.3\% & \textbf{40} & 151\% & 4.83 & 15.7\% & choppy\\
PLTR & 174\% & 2.52 & 25.9\% & \textbf{40} & 195\% & 2.88 & 24.2\% & choppy\\
\rowcolor{bcact} RKLB & 587\% & 3.55 & 34.5\% & \textbf{60} & \textbf{732\%} & \textbf{5.10} & \textbf{26.7\%} & momentum\\
SOFI & 231\% & 4.34 & 15.3\% & \textbf{40} & 241\% & 4.58 & 15.3\% & mild\\
\rowcolor{bckeep} SPOT & 149\% & 6.20 & 5.6\% & \textbf{50} & 149\% & 6.20 & 5.6\% & trender (inert)\\
\bottomrule
\end{tabular}
\caption*{\footnotesize Shaded blue = keep at 50 (flat/inert zone; shortening would only hurt).
Shaded gold = RKLB, the high-impact lengthening. TSLA is shown at 40 for consistency; a 20-day
stop holds Sharpe at 2.71 while cutting drawdown to 17.3\%, if drawdown control is the priority.
Suggested values are round numbers inside each flat zone, not in-sample optima.}
\end{table}

\textbf{Simplest defensible change set.} Move \textbf{TSM, TSLA, VRT, AVGO, PLTR, SOFI to 40},
lengthen \textbf{RKLB to 60}, and \textbf{keep NVDA, VST, SPOT at 50}. Equivalently: a uniform
$\sim$40-day stop for the book, with RKLB the single name held longer.

% ---------------- Robustness ----------------
\section{Does tuning the stop survive out of sample?}

Parameter re-optimisation on this model is a known in-sample mirage --- it does not survive
walk-forward. The stop period is different: it is one low-dimensional, economically-meaningful
knob (holding-period tolerance), and such knobs generalise where a full refit does not.

A split-sample test confirms this. For each name we pick the best-Sharpe \texttt{T2} on the
first half of the sample, apply it \emph{blind} to the second half, and compare its
second-half return to the frozen 50:

\begin{table}[h]\centering\small
\renewcommand{\arraystretch}{1.15}
\begin{tabular}{lcc}
\toprule
\textbf{Result (10 names)} & \textbf{tuned wins} & \textbf{frozen-50 wins} \\
\midrule
Out-of-sample (2nd half) & \textbf{6} (TSM, VRT, AVGO, PLTR, RKLB, SOFI)
                         & 4 (NVDA, TSLA, VST, SPOT) \\
\bottomrule
\end{tabular}
\caption*{\footnotesize 6/10 beats a coin flip and vastly beats full re-optimisation, which
loses out of sample. But TSLA tuned to 40 on the first half did \emph{worse} on the second ---
evidence that the \emph{direction} (choppy $\to$ tighter, momentum $\to$ looser) is robust while
the exact decimal is not.}
\end{table}

\textbf{Practical rule.} Set the stop by regime, not by fitting. Trending / inert-stop names
(NVDA, VST, SPOT) keep the long backstop at 50; choppy names (TSLA, PLTR, AVGO, and mildly TSM,
VRT, SOFI) tighten to $\sim$40; the one pure-momentum outlier (RKLB) lengthens to 60. Adopt the
RKLB change first --- it is the largest and cleanest, improving return, Sharpe and drawdown
together.

\vfill
{\footnotesize\color{bcgrey} Data vintage: current ten-name workbook (\texttt{\_nopq} PLTR
build). Confirm this is the live master before editing cell~T2 on each Model sheet. All results
reproduce the workbook's cached figures to the dollar at T2\,=\,50.}

\end{document}
